% ====================================================================
\section{Completeness results}\label{sec:completeness}
% ====================================================================

\subsection{Membership in P, including binary-encoded times}\label{sec:membership}

The following is the classical bounded-transient argument for symmetric majority networks
\cite{GO80,PS83,GFP85}. We include the proof because we need the explicit transient bound for
binary-encoded times, in the exact model of \cref{def:majority}. It applies to every explicit
instance, not only to the constructed trees.

\begin{proposition}[Membership]\label{prop:membership}
Let $T$ be a finite simple undirected graph with $n$ vertices, and let $W=2|E|+n$. For every initial
configuration, $x(t+2)=x(t)$ for all $t\ge W$. Consequently, $\Pred$ and $\Ever$ are in $\Ptime$, even
when $\tau$ is a large binary integer.
\end{proposition}

\begin{proof}
Use spins $s_v=2x_v-1\in\{-1,1\}$. Put $w_{vu}=w_{uv}=1$ for each edge, $w_{vv}=1$ when $\deg(v)$ is
even (including degree zero), and all other weights $0$. No physical self-loop is added; the diagonal
is used only for the analysis. Every row sum $\sum_u w_{vu}$ is odd, and
\begin{equation}\label{eq:sign}
s_v(t+1)=\operatorname{sign}\Bigl(\sum_u w_{vu}s_u(t)\Bigr),
\end{equation}
because at even degree the diagonal term breaks exactly the ties, and at odd degree there are no
ties. Define
\[
\Phi(t)=\sum_{v,u}w_{vu}\,s_v(t+1)\,s_u(t),\qquad |\Phi(t)|\le\sum_{v,u}w_{vu}=W.
\]
By the symmetry of $w$,
\begin{equation}\label{eq:energy}
\Phi(t+1)-\Phi(t)=\sum_v\bigl(s_v(t+2)-s_v(t)\bigr)\sum_u w_{vu}s_u(t+1).
\end{equation}
The inner sum is an odd integer with sign $s_v(t+2)$ by \eqref{eq:sign}. So each term is $0$ if
$s_v(t+2)=s_v(t)$ and at least $2$ otherwise. Let $t^\ast$ be the first $t$ with $x(t+2)=x(t)$. By
determinism, $x(t+2)=x(t)$ for all $t\ge t^\ast$, and for $t<t^\ast$, $\Phi$ increases by at least $2$
at each step. Since $\Phi$ takes values in $[-W,W]$, we get $2t^\ast\le 2W$.

Algorithm for $\Pred$: simulate up to time $W+1$, which is at most $3n$ for a tree. If $\tau\le W+1$, answer $x_{v^\ast}(\tau)$;
otherwise answer $x_{v^\ast}(W+((\tau-W)\bmod 2))$, using only the parity of $\tau$. Algorithm for
$\Ever$: the states after $W$ repeat with period dividing two, so $z$ is ever $1$ if and only if
$x_z(t)=1$ for some $1\le t\le W+1$. Both run in time polynomial in the explicit input length.
\end{proof}

\subsection{Prediction}

\begin{proof}[Proof of \cref{thmB}]
Membership: \cref{prop:membership}. Hardness: given an instance of $\OMCVP$, which is
$\Ptime$-complete by \cref{thm:mcvp}, \cref{thm:simulation} produces in logspace an instance
$(T,x(0),v^\ast,\tau)$ of $\Pred$ with degrees in $\{1,3,4\}$, polynomial $\tau$, and
$x_{v^\ast}(\tau)=1$ if and only if the circuit outputs $1$. Composition of logspace reductions
\cite{Jon75,LL76} completes the proof.
\end{proof}

\subsection{Eventual activation}\label{sec:ever}

The clock theorem implements the processor controls only up to a finite horizon. After $\Hz$ the
generators alternate, and the processor may behave arbitrarily (\cref{rem:afterH}). Prediction at the
single time $\tau$ is unaffected, but eventual activation is not: a false instance must never
activate the designated vertex, even after $\Hz$. We need a reader that rejects in the false case
whatever happens outside its window.

\begin{lemma}[Finite-window reader]\label{lem:reader}
Let $q$ be a vertex with an arbitrary history, except that $x_q(T_0)=x_q(T_0+1)=d$. Attach a new
vertex $z$ adjacent to $q$, to a private leaf $\ell$, and to the roots of two private copies of
$P_{T_0,T_0+2}$ (\cref{lem:pulses}). Initialise $z=\ell=0$. Then $z$ has degree four, $\ell$ has
degree one, the pulses use degrees $1$ and $3$, and the edge $qz$ counts towards the degree of $q$.
Moreover:
\begin{enumerate}[label=(\alph*)]
\item $x_z(t)=0$ for all $t\le T_0$, whatever the history of $q$;
\item $x_z(t)=1$ for some $t$ if and only if $d=1$, in which case $x_z(T_0+1)=1$.
\end{enumerate}
\end{lemma}

\begin{proof}
The pulses are robust and their consumer is $z$, so both are $1$ exactly at times $T_0$ and $T_0+1$,
and $0$ at all other times, \emph{including all times after $T_0+1$}: their tail is exactly zero for
every $t\ge T_0+2$. Before $T_0$, if $z=\ell=0$ then only $q$ can vote one at $z$, so $z$ stays $0$,
and the leaf copies $z$. Hence $z=\ell=0$ up to time $T_0$. At control time $T_0$,
$x_z(T_0+1)=\MAJ_5(0,d,0,1,1)=d$ and $x_\ell(T_0+1)=0$. At $T_0+1$,
$x_z(T_0+2)=\MAJ_5(d,d,0,1,1)=d$ and $x_\ell(T_0+2)=d$. From control time $T_0+2$ on, both pulses
are $0$. If $d=0$, the invariant $z=\ell=0$ is restored at $T_0+2$ and persists for every later
history of $q$, since only $q$ can vote one. If $d=1$, the witness $x_z(T_0+1)=1$ has occurred;
retaining $1$ forever is not needed.
\end{proof}

The reader does not require the input to start at zero, nor that it was never $1$ earlier. Nor does
it require the phase-defect clocks to keep \HOLD\ forever.

\begin{construction}[The activation tree]\label{con:ever}
Let $\mathcal P$ be in normal form; its last instruction is $R_m\leftarrow R_o\wedge R_o$, of
background $0$. Modify \cref{con:tree} as follows.
\begin{enumerate}[label=(\arabic*)]
\item \emph{Reader at the far port.} Remove the far clock $\kappa_m$, and attach the reader vertex $z$
  of \cref{lem:reader} to $q=\ep{R_m}{\kappa_m}$, in the same data port. This adds no edge at the
  register endpoint, which still has degree four.
\item \emph{Final instruction.} In $\Sigma_m$, omit $\prim{SET}(R_m\mathbin{:=}b)$ from
  \prim{PREPARE}, so that $R_m$ stays $00$ and only $C_m=K_m=00$ are prepared. Replace phases 11
  and 12 by the \emph{final write}:
  \begin{align*}
  &\prim{ROUTE}(r_m;E_{m-1}\to C_m;E_m) &&\text{during }[s,s+6),\\
  &\prim{MOVE}_0(C_m\to K_m) &&\text{during }[s+6,s+9),\\
  &\prim{MOVE}_0(K_m\to R_m) &&\text{during }[s+9,s+12),
  \end{align*}
  starting at the end $s$ of phase 10, which is empty since $k=m$. Let $\tau_E=s+12$, and keep the
  rest schedule of background $0$ afterwards.
\item \emph{Window.} Use $T_0=\tau_E$ in the reader, and the horizon $\Hz=\tau_E+3$ for all
  generators.
\end{enumerate}
The designated vertex is $z$, with $x_z(0)=0$.
\end{construction}

\begin{lemma}[Final write]\label{lem:finalwrite}
In the activation tree, the processor vertices follow the forced evolution of the modified schedule
up to time $\Hz+1$. At times $\tau_E$ and $\tau_E+1$, $R_m=dd$ with $d=\val(\mathcal P)$. All blocks
and control events of the modified schedule are admissible.
\end{lemma}

\begin{proof}
\emph{Before the final instruction.} The port of $\ep{R_m}{\kappa_m}$ now carries $x_z$ instead of a
far clock. By \cref{lem:reader}(a), $x_z(t)=0$ for $t\le\tau_E$, for every history of $R_m$. The cell
$R_m$ is in \HOLD\ and uniform $00$ from time $0$ until the last move. It is never read, because every
operand index is smaller than $m$. It is not prepared, and normalisation leaves registers in \HOLD.
So the content of its far port is never used before the final write, and
\cref{lem:prepare,lem:fetch,lem:normalize,prop:instruction} hold verbatim for all instructions
$k<m$. Within $\Sigma_m$, phases 1 to 10 are as before, except that $R_m=00$ instead of $bb$; these
phases never involve $R_m$. No read enters the branch of $R_m$.

\emph{Final write ($c=0$, $b=1$).} At $s$ we have $E_{m-1}=dd$, $E_m=11$ (the pilot), $C_m=K_m=00$,
$R_m=00$ and $x_{r_m}=0$. The route has destination $C_m=00$, whose far port is $\ep{K_m}{C_m}=0$, and
bias $E_m=11$. Its source's far port is $r_{m-1}$, whose arms are $E_{m-2}=0$, $C_{m-1}=1$ and
$E_{m-1}$, so it stays $0$. By \cref{lem:route}, $C_m=dd$ and $E_{m-1}=00$. Then
$\prim{MOVE}_0(C_m\to K_m)$: the destination is $K_m=00$ with far port $\ep{R_m}{K_m}=0$, which gives
$C_m=00$ and $K_m=dd$. Then $\prim{MOVE}_0(K_m\to R_m)$: the destination is $R_m=00$, and its far
port is $x_z$, which is $0$ at control times $\tau_E-3,\tau_E-2,\tau_E-1\le\tau_E$. This gives
$K_m=00$ and $R_m=dd$ at $\tau_E$. During these moves $r_m$ keeps $0$, because its self-vote, $h_m$
and $E_{m-1}=0$ are $0$. Afterwards $R_m$ is in \HOLD\ and uniform, so it holds $dd$ at $\tau_E+1$,
whatever $z$ does. Since this is the final write, neither restoring $C_m,K_m$ nor returning the whole
processor to a normal configuration is needed.

\emph{Admissibility.} The route at $r_m$ has reference parity $c=0$, and phase 10 ends at parity $0$
(\cref{lem:sigmatiming}). $\prim{MOVE}_0(C_m\to K_m)$ is the event $(0,0,6)$, and
$\prim{MOVE}_0(K_m\to R_m)$ is $(1,0,9)$. Both give $0$ by \eqref{eq:startparity}. The omitted
$\prim{SET}(R_m\mathbin{:=}b)$ only removes a block. The conclusions of \cref{thm:forced}\ref{forced:phase}
and of \cref{prop:correct} therefore hold for the modified schedule, with $\Hz=\tau_E+3$.
\end{proof}

\begin{proof}[Proof of \cref{thmC} and of the activation part of \cref{thm:simulation}]
\emph{Topology and size.} The activation tree is obtained from a tree with the far clock $\kappa_m$
removed by attaching the reader through the single edge $qz$. The reader is itself a tree, with
$z$ of degree four, $\ell$ of degree one, and pulses with degrees $1$ and $3$. So it is a tree with
degrees in $\{1,3,4\}$. The two pulses have size $O((\tau_E+3)^2)$.

\emph{Correctness.} By \cref{lem:finalwrite}, $x_q(\tau_E)=x_q(\tau_E+1)=d=\val(\mathcal P)$ in the
closed tree. By \cref{lem:reader}, $z$ is ever $1$ if and only if $d=1$. If the program outputs $0$,
then $z$ stays $0$ forever, even though the phase-defect clocks alternate after $\Hz$ and may change
the processor arbitrarily.

\emph{Uniformity.} The reader and its pulses are emitted in logspace (\cref{lem:genlogspace}, or
directly from the explicit description of \cref{lem:pulses}), and the modified schedule is
interpreted as in \cref{prop:logspace}.

\emph{Completeness.} Membership is \cref{prop:membership}, and hardness follows from \cref{thm:mcvp}
as in the proof of \cref{thmB}.
\end{proof}

\subsection{Scope}

\begin{remark}\label{rem:scope}
\begin{itemize}
\item The construction proves hardness already at treewidth one.
\item It does not show that four is the smallest possible maximum degree. The classification at
  maximum degree three remains unresolved in the source notes.
\item The construction of degree five with only odd degrees is a separate, earlier result.
\item Nothing here separates $\Ptime$ from $\NC$, and no claim of bibliographic priority or
  external validation is made.
\end{itemize}
\end{remark}
